% GENERATED FILE. Edit canonical lessons, then run npm run generate:books.
% canonical-source-hash: fnv1a64:1bb95646b7a5da6a
% canonical-lessons: BN-C451-ghinghine, BN-C451-bitrishna, BN-C451-poroshrikator, BN-C451-ghrina-kora, BN-C451-nak-sitkano

\chapter{Disgust and Envy}
\label{ch:bn-disgust-and-envy}
\hlchaptermodality{451}

\begin{chapteropening}
\textbf{By the end of this chapter:} \emph{I can say that I hate something or find it disgusting.}

Complete the last lesson of chapter 451: I can say that I hate something or find it disgusting.
\end{chapteropening}

\section[\texorpdfstring{\bn{ঘিনঘিনে} (ghinghine)}{ঘিনঘিনে (ghinghine)}]{\bn{ঘিনঘিনে} (ghinghine) — disgusting; revolting}

\label{lesson:BN-C451-ghinghine}



\begin{quote}
Before the new one: say the Bengali for to taunt, then the Bengali for to belittle.
\end{quote}

\subsection*{You'll want to know: \bn{ঘিনঘিনে}}
\textbf{\bn{ঘিনঘিনে}} — \emph{ghinghine} — \textquotedblleft{}disgusting; revolting\textquotedblright{}.

\begin{quote}
\bn{আমার} \bn{কাছে} \bn{মাকড়সা} \bn{ঘিনঘিনে} \bn{লাগে।}
\end{quote}

\begin{quote}
\emph{āmār kāchhe mākôṛsā ghinghine lāge.}
\end{quote}

\begin{quote}
I find spiders revolting.
\end{quote}

\begin{cousinweb}
\textbf{\bn{ঘিনঘিনে}} (\emph{ghinghine}) doubles \emph{ghin}, \textquotedblleft{}loathing\textquotedblright{}, which CDIAL lists for Bengali under Sanskrit \emph{ghṛṇā} (CDIAL 4500), and adds the adjective ending \emph{-e}.
\end{cousinweb}

\begin{grammarlens}[title={What you've built}]
One more word for describing what you lived and felt.
\end{grammarlens}

\subsection*{Your turn}
\begin{itemize}
\raggedright
  \item \emph{Say it:} \emph{ghinghine}
  \item \emph{Say it:} \emph{ghinghine}, once more
  \item \emph{Say it:} \emph{āmār kāchhe mākôṛsā ghinghine lāge.} — I find spiders revolting.
  \item \emph{Recall:} say the Bengali for to taunt, then the Bengali for to belittle, then say \emph{ghinghine} again
\end{itemize}

\begin{tcolorbox}[breakable,colback=teal!4,colframe=teal!35!black,title={Before you move on}]
What does \bn{ঘিনঘিনে} mean? (\textquotedblleft{}Disgusting; revolting\textquotedblright{}.) Say it once more.
\end{tcolorbox}
\section[\texorpdfstring{\bn{বিতৃষ্ণা} (bitrishnā)}{বিতৃষ্ণা (bitrishnā)}]{\bn{বিতৃষ্ণা} (bitrishnā) — aversion; distaste}

\label{lesson:BN-C451-bitrishna}



\begin{quote}
Before the new one: say the Bengali for to belittle, then the Bengali for revolting.
\end{quote}

\subsection*{You'll want to know: \bn{বিতৃষ্ণা}}
\textbf{\bn{বিতৃষ্ণা}} — \emph{bitrishnā} — \textquotedblleft{}aversion; distaste\textquotedblright{}.

\begin{quote}
\bn{রোজ} \bn{একই} \bn{খাবার} \bn{খেয়ে} \bn{অরুণের} \bn{বিতৃষ্ণা} \bn{এল।}
\end{quote}

\begin{quote}
\emph{roj eki khābār kheye ôruner bitrishnā elô.}
\end{quote}

\begin{quote}
Eating the same food every day, Arun grew sick of it.
\end{quote}

\begin{cousinweb}
\textbf{\bn{বিতৃষ্ণা}} (\emph{bitrishnā}) is a tatsama, Sanskrit \emph{vitṛṣṇā}, which Monier-Williams gives as freedom from desire, \emph{vi-} \textquotedblleft{}without\textquotedblright{} on \emph{tṛṣṇā} \textquotedblleft{}thirst\textquotedblright{}; Bengali has pushed it on to distaste, a sense the dictionary does not give.
\end{cousinweb}

\begin{grammarlens}[title={What you've built}]
One more word for describing what you lived and felt.
\end{grammarlens}

\subsection*{Your turn}
\begin{itemize}
\raggedright
  \item \emph{Say it:} \emph{bitrishnā}
  \item \emph{Say it:} \emph{bitrishnā}, once more
  \item \emph{Say it:} \emph{roj eki khābār kheye ôruner bitrishnā elô.} — Eating the same food every day, Arun grew sick of it.
  \item \emph{Recall:} say the Bengali for to belittle, then the Bengali for revolting, then say \emph{bitrishnā} again
\end{itemize}

\begin{tcolorbox}[breakable,colback=teal!4,colframe=teal!35!black,title={Before you move on}]
What does \bn{বিতৃষ্ণা} mean? (\textquotedblleft{}Aversion; distaste\textquotedblright{}.) Say it once more.
\end{tcolorbox}
\section[\texorpdfstring{\bn{পরশ্রীকাতর} (pôrôshrīkātôr)}{পরশ্রীকাতর (pôrôshrīkātôr)}]{\bn{পরশ্রীকাতর} (pôrôshrīkātôr) — envious of others' good fortune}

\label{lesson:BN-C451-poroshrikator}



\begin{quote}
Before the new one: say the Bengali for revolting, then the Bengali for aversion.
\end{quote}

\subsection*{You'll want to know: \bn{পরশ্রীকাতর}}
\textbf{\bn{পরশ্রীকাতর}} — \emph{pôrôshrīkātôr} — \textquotedblleft{}envious of others' good fortune\textquotedblright{}.

\begin{quote}
\bn{প্রতিবেশীর} \bn{নতুন} \bn{গাড়ি} \bn{দেখে} \bn{লোকটা} \bn{পরশ্রীকাতর} \bn{হল।}
\end{quote}

\begin{quote}
\emph{prôtibeshīr nôtun gāṛi dekhe lokṭā pôrôshrīkātôr hôlô.}
\end{quote}

\begin{quote}
The man grew envious when he saw his neighbour's new car.
\end{quote}

\begin{cousinweb}
\textbf{\bn{পরশ্রীকাতর}} (\emph{pôrôshrīkātôr}) puts together Sanskrit \emph{para}, \textquotedblleft{}another\textquotedblright{}, \emph{śrī}, \textquotedblleft{}good fortune\textquotedblright{}, and \emph{kātara}, \textquotedblleft{}despairing, agitated\textquotedblright{} (Monier-Williams): pained by someone else's good luck.
\end{cousinweb}

\begin{grammarlens}[title={What you've built}]
One more word for describing what you lived and felt.
\end{grammarlens}

\subsection*{Your turn}
\begin{itemize}
\raggedright
  \item \emph{Say it:} \emph{pôrôshrīkātôr}
  \item \emph{Say it:} \emph{pôrôshrīkātôr}, once more
  \item \emph{Say it:} \emph{prôtibeshīr nôtun gāṛi dekhe lokṭā pôrôshrīkātôr hôlô.} — The man grew envious when he saw his neighbour's new car.
  \item \emph{Recall:} say the Bengali for revolting, then the Bengali for aversion, then say \emph{pôrôshrīkātôr} again
\end{itemize}

\begin{tcolorbox}[breakable,colback=teal!4,colframe=teal!35!black,title={Before you move on}]
What does \bn{পরশ্রীকাতর} mean? (\textquotedblleft{}Envious of others' good fortune\textquotedblright{}.) Say it once more.
\end{tcolorbox}
\section[\texorpdfstring{\bn{ঘৃণা} \bn{করা} (ghrinā kôrā)}{ঘৃণা করা (ghrinā kôrā)}]{\bn{ঘৃণা} \bn{করা} (ghrinā kôrā) — to hate; to despise}

\label{lesson:BN-C451-ghrina-kora}



\begin{quote}
Before the new one: say the Bengali for aversion, then the Bengali for envious of others' good fortune.
\end{quote}

\subsection*{You'll want to know: \bn{ঘৃণা} \bn{করা}}
\textbf{\bn{ঘৃণা} \bn{করা}} — \emph{ghrinā kôrā} — \textquotedblleft{}to hate; to despise\textquotedblright{}.

\begin{quote}
\bn{মীরা} \bn{মিথ্যেকে} \bn{ঘৃণা} \bn{করে।}
\end{quote}

\begin{quote}
\emph{mīrā mittheke ghrinā kôre.}
\end{quote}

\begin{quote}
Mira hates lies.
\end{quote}

Past (he/she …): \bn{তিনি} \bn{যুদ্ধকে} \bn{ঘৃণা} \bn{করলেন}, \emph{tini juddhôke ghrinā kôrlen}.

\begin{cousinweb}
\textbf{\bn{ঘৃণা} \bn{করা}} (\emph{ghrinā kôrā}) is Sanskrit \emph{ghṛṇā}, which Monier-Williams gives for warm feeling and pity and also for contempt and disgust; Bengali keeps the loathing, with \emph{kôrā}, \textquotedblleft{}to do\textquotedblright{} (CDIAL 2814).
\end{cousinweb}

\begin{grammarlens}[title={What you've built}]
One more word for describing what you lived and felt.
\end{grammarlens}

\subsection*{Your turn}
\begin{itemize}
\raggedright
  \item \emph{Say it:} \emph{ghrinā kôrā}
  \item \emph{Say it:} \emph{ghrinā kôrā}, once more
  \item \emph{Say it:} \emph{mīrā mittheke ghrinā kôre.} — Mira hates lies.
  \item \emph{Recall:} say the Bengali for aversion, then the Bengali for envious of others' good fortune, then say \emph{ghrinā kôrā} again
\end{itemize}

\begin{tcolorbox}[breakable,colback=teal!4,colframe=teal!35!black,title={Before you move on}]
What does \bn{ঘৃণা} \bn{করা} mean? (\textquotedblleft{}To hate; to despise\textquotedblright{}.) Say it once more.
\end{tcolorbox}
\section[\texorpdfstring{\bn{নাক} \bn{সিঁটকানো} (nāk sĩṭkāno)}{নাক সিঁটকানো (nāk sĩṭkāno)}]{\bn{নাক} \bn{সিঁটকানো} (nāk sĩṭkāno) — to turn up one's nose}

\label{lesson:BN-C451-nak-sitkano}



\begin{quote}
Before the new one: say the Bengali for envious of others' good fortune, then the Bengali for to hate.
\end{quote}

\subsection*{You'll want to know: \bn{নাক} \bn{সিঁটকানো}}
\textbf{\bn{নাক} \bn{সিঁটকানো}} — \emph{nāk sĩṭkāno} — \textquotedblleft{}to turn up one's nose\textquotedblright{}.

\begin{quote}
\bn{সবজি} \bn{দেখে} \bn{ছেলেটা} \bn{নাক} \bn{সিঁটকাল।}
\end{quote}

\begin{quote}
\emph{sôbji dekhe chheleṭā nāk sĩṭkālô.}
\end{quote}

\begin{quote}
The boy turned up his nose at the vegetables.
\end{quote}

Past (he/she …): \bn{তিনি} \bn{খাবার} \bn{দেখে} \bn{নাক} \bn{সিঁটকালেন}, \emph{tini khābār dekhe nāk sĩṭkālen}.

\begin{cousinweb}
\textbf{\bn{নাক} \bn{সিঁটকানো}} (\emph{nāk sĩṭkāno}) is \bn{নাক} (taught), the nose (CDIAL 6909), with \emph{sĩṭkāno}, \textquotedblleft{}to wrinkle up\textquotedblright{}, a verb whose origin is not settled.
\end{cousinweb}

\begin{grammarlens}[title={What you've built}]
One more word for describing what you lived and felt.
\end{grammarlens}

\subsection*{Your turn}
\begin{itemize}
\raggedright
  \item \emph{Say it:} \emph{nāk sĩṭkāno}
  \item \emph{Say it:} \emph{nāk sĩṭkāno}, once more
  \item \emph{Say it:} \emph{sôbji dekhe chheleṭā nāk sĩṭkālô.} — The boy turned up his nose at the vegetables.
  \item \emph{Recall:} say the Bengali for envious of others' good fortune, then the Bengali for to hate, then say \emph{nāk sĩṭkāno} again
\end{itemize}

\begin{tcolorbox}[breakable,colback=teal!4,colframe=teal!35!black,title={Before you move on}]
What does \bn{নাক} \bn{সিঁটকানো} mean? (\textquotedblleft{}To turn up one's nose\textquotedblright{}.) Say it once more.
\end{tcolorbox}
